\documentclass{beamer}
\usetheme{Madrid}
\usecolortheme{default}
\usepackage{listings}
\usepackage{graphicx}
\usepackage{booktabs}
\usepackage{amsmath}
\lstset{basicstyle=\ttfamily\small, breaklines=true, language=Python}

\title{GAMs to Predict TAVG PC Coefficients from Geography}
\author{CSE255}
\institute{}
\date{}

\begin{document}

\begin{frame}
\titlepage
\end{frame}

%==============================================================================
\section{Introduction}
%==============================================================================

\begin{frame}{Overview}
\begin{itemize}
\item Use PCA of yearly TAVG (temperature) profiles
\item Predict PC1, PC2, PC3 from latitude, distance to coast, longitude, elevation
\item Start with visualizing the PCA structure, then fit GAMs
\end{itemize}
\end{frame}

%==============================================================================
\section{PCA Visualization}
%==============================================================================

\begin{frame}{Load PCA Results}
Load pre-computed TAVG PCA from pickle:
\begin{lstlisting}
with open(PCA_PATH, 'rb') as f:
    pca_results = pickle.load(f)
pca_mean = pca_results['mean']        # shape (365,)
pca_components = pca_results['components']  # shape (10, 365)
evr = pca_results['explained_variance_ratios']
\end{lstlisting}
\vspace{0.3em}
Top 3 variance ratios: PC1 $\approx$ 76\%, PC2 $\approx$ 17\%, PC3 $\approx$ 3\%
\end{frame}

\begin{frame}{Mean and Top 3 Eigenvectors}
\begin{itemize}
\item \textbf{Left y-axis:} mean TAVG profile (day 1--365)
\item \textbf{Right y-axis (twin):} PC1, PC2, PC3 eigenvectors
\item PC1 captures main seasonal cycle; PC2/PC3 capture phase and higher-order patterns
\end{itemize}
\vspace{0.3em}
% Export from notebook and uncomment:
% \includegraphics[width=0.85\textwidth]{mean_and_pcs.pdf}
\end{frame}

\begin{frame}{Variance Explained}
\begin{itemize}
\item Bar chart: percentage variance explained per PC (PC1--PC10)
\item Cumulative line overlay (e.g., red)
\item Top 3 PCs explain $\approx$ 95\% of variance for TAVG
\end{itemize}
\vspace{0.3em}
% \includegraphics[width=0.85\textwidth]{variance_explained.pdf}
\end{frame}

%==============================================================================
\section{Data Preparation}
%==============================================================================

\begin{frame}{Compute PC Coefficients}
Project each station-year's daily TAVG onto PCA components:
\begin{lstlisting}
daily_values = np.nan_to_num(tavg[daily_cols].values, nan=0.0)
normalized = daily_values - pca_mean[np.newaxis, :]
for i in range(3):
    tavg[f'PC{i+1}'] = normalized @ pca_components[i]
\end{lstlisting}
PC coefficient = (data $-$ mean) $\cdot$ component
\end{frame}

\begin{frame}{Aggregate by Station}
\begin{lstlisting}
agg = tavg.groupby('ID').agg({
    'latitude': 'first', 'longitude': 'first', 'elevation': 'first',
    'dist_to_coast': 'first', 'PC1': 'mean', 'PC2': 'mean', 'PC3': 'mean'
}).reset_index()
\end{lstlisting}
\begin{itemize}
\item One row per station
\item Mean PC coefficients across years
\item Drop invalid elevation ($-999.9$)
\end{itemize}
\end{frame}

\begin{frame}{Predictor Matrix}
\begin{lstlisting}
predictor_cols = ['latitude', 'dist_to_coast', 'longitude', 'elevation']
X = agg[predictor_cols].to_numpy(dtype=float)
\end{lstlisting}
Optional: subsample to 50k rows if large. Response: $y \in \{\mathrm{PC1}, \mathrm{PC2}, \mathrm{PC3}\}$.
\end{frame}

%==============================================================================
\section{GAM Modeling}
%==============================================================================

\begin{frame}{GAM Specification}
For each response PC1, PC2, PC3:
\begin{lstlisting}
gam = LinearGAM(s(0, n_splines=20) + l(1) + l(2) + l(3))
gam.gridsearch(X, y, progress=False)
\end{lstlisting}
\begin{itemize}
\item \texttt{s(0)}: spline on latitude
\item \texttt{l(1)}, \texttt{l(2)}, \texttt{l(3)}: linear terms for dist\_to\_coast, longitude, elevation
\end{itemize}
\end{frame}

\begin{frame}{Partial Dependence Plots}
\begin{itemize}
\item 3$\times$4 grid: rows = PC1, PC2, PC3; columns = predictors
\item For each term: \texttt{generate\_X\_grid}, \texttt{partial\_dependence(..., width=0.95)}
\item Shaded bands show 95\% confidence intervals
\end{itemize}
\vspace{0.3em}
% \includegraphics[width=0.9\textwidth]{partial_dependence.pdf}
\end{frame}

%==============================================================================
\section{Goodness of Fit}
%==============================================================================

\begin{frame}{Metrics Table}
\begin{center}
\begin{tabular}{lcccc}
\toprule
\textbf{Model} & \textbf{R²} & \textbf{MAE} & \textbf{RMSE} & \textbf{AIC} \\
\midrule
PC1 & 0.77 & 531 & 739 & 265653 \\
PC2 & 0.57 & 333 & 468 & 248486 \\
PC3 & 0.33 & 88 & 120 & 197497 \\
\bottomrule
\end{tabular}
\end{center}
\end{frame}

\begin{frame}{Column Definitions}
\begin{itemize}
\item \textbf{Model}: PC coefficient predicted
\item \textbf{R²}: Pseudo R² (deviance explained); higher = better fit
\item \textbf{MAE}: Mean absolute error; lower = better
\item \textbf{RMSE}: Root mean squared error; lower = better
\item \textbf{AIC}: Akaike Information Criterion; lower = better (balances fit vs.\ complexity)
\end{itemize}
\end{frame}

\begin{frame}{Akaike Information Criterion (AIC)}
\textbf{Formula:} AIC $= 2k - 2\ln(\hat{L})$
\begin{itemize}
\item $k$ = number of parameters
\item $\hat{L}$ = maximized likelihood
\item Equivalently: AIC $= D + 2k$, with deviance $D = -2\ln(\hat{L})$
\end{itemize}
\vspace{0.5em}
\textbf{Interpretation:}
\begin{itemize}
\item $-2\ln(\hat{L})$ measures lack of fit
\item $2k$ penalizes complexity
\item Lower AIC is better; $\Delta$AIC $> 2$ often meaningful, $> 6$ strong preference
\end{itemize}
\end{frame}

\begin{frame}{Summary}
\begin{itemize}
\item TAVG PCA: mean + top 3 eigenvectors, variance explained
\item Geographic predictors: latitude (spline), dist\_to\_coast, longitude, elevation (linear)
\item LinearGAM fits PC1, PC2, PC3; PC1 best explained (R² $\approx$ 0.77)
\item Partial dependence plots reveal non-linear effects of latitude
\end{itemize}
\end{frame}

\end{document}
